\documentclass[11pt]{article}
\usepackage[margin=1in]{geometry}
\usepackage{amsmath,amssymb}
\usepackage{enumitem}
\usepackage{array,longtable}
\usepackage[hidelinks]{hyperref}
\setlist[enumerate]{itemsep=0.45em,topsep=0.5em,parsep=0pt}

\title{Response to the Referees\\
\large Manuscript ITA-2026-0032}
\author{}
\date{12 July 2026}

\begin{document}
\maketitle

\noindent
We thank the referees for the detailed reports.  The revised manuscript has
the title \emph{Canonical Zeckendorf Normalization and Sharp Iteration Depth
of the Berstel Adder}.  It has been substantially rewritten to clarify the contribution, to align the
terminology with the standard LSDF/MSDF and subsequential-transducer
literature, and to avoid claiming as new results that are already classical in
the work of Frougny, Sakarovitch, Berstel, and subsequent automata-theoretic
verification papers.  The main revision is that the paper now foregrounds
exact invariants for the canonical LSDF normalizer: infinite coordinate
anticipation on the genuine addition image, nonextendability at alternating
one-sided limit tails, maximal online commitment deficit $B(n)=n$, maximal
prefix erasure $\Delta(n)=\Delta_{\operatorname{Add}_2}(n)=n$, and linear
input-scale reduced-residual growth $R_{\mathrm{in}}(n)\ge n$.  The
non-subsequentiality statement is
now presented as a consequence of these sharper results.  In addition, the
revision clarifies the Berstel state model: the classical ten-state table is
a raw complete-output kernel, while standard output-delay normalization gives
a six-state subsequential quotient, and the six reduced residuals are
pairwise distinct.  We have also added a new section on iterating the
Berstel recoding to the greedy Zeckendorf normal form, proving the sharp
binary cleanup bound $\lfloor L/2\rfloor$ and the sharp worst-case depth
$\lceil n/2\rceil$ on genuine addition inputs through a cascade-depth
invariant.  The extremal genuine-addition bound is now attained by trimmed
inputs, so it does not rely on leading-zero padding.

For ease of verification, Table~\ref{tab:response-comparison} identifies the
principal current statements, their locations in the revised manuscript, and
the closest prior results.  We use cautious novelty language throughout: to
the best of our knowledge, the exact quantitative formulas identified in the
last column are not stated in the cited sources.

\small
\begin{longtable}{>{\raggedright\arraybackslash}p{0.14\textwidth}
                  >{\raggedright\arraybackslash}p{0.12\textwidth}
                  >{\raggedright\arraybackslash}p{0.28\textwidth}
                  >{\raggedright\arraybackslash}p{0.34\textwidth}}
\caption{Theorem-by-theorem comparison with the closest prior results.}
\label{tab:response-comparison}\\
\hline
Current result & Location & Closest prior result & Precise increment claimed in the revision\\
\hline
\endfirsthead
\multicolumn{4}{l}{\small Table~\thetable\ continued from the preceding page.}\\
\hline
Current result & Location & Closest prior result & Precise increment claimed in the revision\\
\hline
\endhead
Theorem 4.3
& p.~8
& Frougny, \emph{Representations of numbers and finite automata} (1992)
[revised Ref.~7], and Sakarovitch, \emph{Easy multiplications I} (1987)
[revised Ref.~14]: directional transducer background for Fibonacci
normalization.
& For canonical greedy LSDF output, every fixed coordinate has infinite
anticipation, both on all ternary finite-support inputs and on the genuine
addition image.\\
Theorem 4.5
& pp.~9--10
& Frougny (1992) [revised Ref.~7] and Sakarovitch (1987) [revised Ref.~14]
provide the directional normalization background; they are not cited for the
present topological formula.
& Explicit alternating one-sided tails at which the canonical coordinate maps
have no continuous extension from finite-support inputs.\\
Theorem 4.6
& p.~11
& The directional-transducer background in Frougny,
\emph{Representations of numbers and finite automata} (1992) [revised Ref.~7],
and Sakarovitch, \emph{Easy multiplications I} (1987) [revised Ref.~14].
& Exact finite-word commitment and prefix-erasure formulas
$B(n)=\Delta(n)=n$.\\
Theorem 4.8
& p.~12
& The directional normalization background in revised Refs.~7 and 14, without
the present restricted-domain formula.
& Exact maximal prefix erasure
$\Delta_{\operatorname{Add}_2}(n)=n$ inside the genuine addition image.\\
Theorem 4.9
& p.~13
& The directional and qualitative finite-state background in revised
Refs.~7 and 14; no finite input-scale residual count is attributed to them.
& A finite input-scale residual count and the explicit lower bound
$R_{\mathrm{in}}(n)\ge n$.\\
Propositions 6.4--6.5
& pp.~17--19
& The classical ten-state Berstel kernel; Mohri (1994), \emph{Minimization of
Sequential Transducers} [revised Ref.~12], for string-to-string sequential
transducer minimization.
& Output-delay normalization gives six pairwise distinct reduced residuals and
the minimal six-state quotient when an initial output word is allowed.\\
Theorem 7.5
& p.~23
& Frougny (1999), Proposition 13 (pp. 98--99), for the base-$\varphi$ online
machine, and Corollary 4 (p. 100), for Fibonacci numeration with terminal
output [revised Ref.~8]; Labb\'e--Lep\v{s}ov\'a (2023), Theorem 3.1 [revised
Ref.~11], for a constructive treatment of the Berstel recoding.
& Exact worst-case iteration depth $\lfloor L/2\rfloor$ for trimmed binary
inputs.\\
Theorem 7.9
& p.~26
& The same classical one-pass recoding results, which do not quantify repeated
cleanup to greedy form.
& Exact worst-case depth $\lceil n/2\rceil$ on genuine addition inputs, attained
without leading-zero padding.\\
\hline
\end{longtable}
\normalsize

\section*{Response to Referee 1}

\begin{enumerate}[leftmargin=*]
\item \textbf{Prior work on right/left subsequential normalization.}
The introduction now explicitly cites and discusses Frougny's finite-automaton
and online-addition work, Sakarovitch's decomposition of Fibonacci
normalization through one-sided subsequential transformations, Frougny's
confluent linear numeration systems, Berstel's adder, and the mechanical
verification work of Mousavi--Schaeffer--Shallit.  The manuscript now states
that it is not claiming a first non-subsequentiality theorem for Fibonacci
normalization.  The revised contribution is the exact package of LSDF
noncausality invariants for the canonical greedy normalizer, including the
genuine-addition restrictions described above.  The introduction now contains
a result-by-result comparison table recording the prior domain, reading
direction, output convention, and the exact increment proved in this paper
(revised manuscript, Table~1, p.~3).  The principal quantitative statements
are Theorems~4.3, 4.5, 4.6, 4.8, and 4.9 (pp.~8--13).
We agree that the qualitative failure of one-direction subsequential
normalization is classical.  Corollary~4.12 is therefore not claimed as new;
our novelty claim is restricted to the exact quantitative invariants and
their genuine-addition-domain versions.

\item \textbf{Berstel adder and the ten-state presentation.}
The title, abstract, introduction, and Sections 5--6 were revised to remove
the impression that the ten-state Berstel presentation is new.  We no longer
claim unrestricted deterministic subsequential state complexity $10$.  The
revised Section 6 distinguishes the displayed ten-state complete-output kernel
from unrestricted subsequential minimization and gives the standard
output-delay quotient, which has six pairwise distinct reduced residuals.  The
minimality statement now cites Mohri's string-to-string minimization result
(1994, \emph{Minimization of Sequential Transducers}, revised Ref.~12) and
also includes a self-contained lower-bound argument: two prefixes reaching
one state induce the same normalized residual, whereas the six reachable
residuals are pairwise distinct (Propositions~6.4--6.5, pp.~17--19).  The
six-state minimality statement uses the standard convention that an initial
output word is allowed; without initial output, one dummy initial state is
added.

\item \textbf{Mousavi--Schaeffer--Shallit presentation.}
We corrected our earlier description.  The referee's count of $16$ excludes
the unique dead state $0$; the published state set is $Q=\{0,\ldots,16\}$,
hence the machine is a $17$-state three-tape DFA with $16$ nondead states.  It
recognizes padded triples $(x,y,z)_F$ with $x+y=z$ and is a relation acceptor,
not an equivalent state realization of the Berstel recoding [revised Ref.~13].
The revised manuscript states this distinction and cites
Labb\'e--Lep\v{s}ov\'a [revised Ref.~11, Thm.~3.1] for a direct constructive
proof of the recoding (discussion preceding Proposition~5.1, pp.~14--15).

\item \textbf{Locality lower bound.}
The local-propagation lower bound has been moved out of the main line of the
paper and placed in an appendix as a witness for the standard obstruction.
The introduction no longer presents it as a main contribution, and the
appendix now explicitly cites Frougny's Proposition~14 [revised Ref.~8,
p.~99] and describes the
statement only as a finite-scale light-cone formulation.  The revised
statement fixes an input scale $m$ and assumes correctness only for inputs
supported in $\{1,\ldots,m\}$ (Theorem~A.1, p.~27).

\item \textbf{Confluent numeration systems.}
The introduction now states that the paper is Fibonacci-specific and does not
claim a new normalization theory beyond that setting (p.~2).
We have not attempted that generalization in the revision: the new principal
result concerns the concrete iteration dynamics of the Berstel recoding, and
extension to other confluent systems is a separate problem requiring different
transducers and carry invariants.
\end{enumerate}

\section*{Response to Referee 2}

\begin{enumerate}[leftmargin=*]
\item \textbf{Contribution relative to Frougny and Sakarovitch.}
The introduction has been rewritten.  It now says precisely which parts are
classical and what is retained as the paper's contribution: a self-contained
canonical-section and rewrite presentation; explicit LSDF witness families
proving the exact infinite anticipation profile
$A_{\{0,1,2\}^{(\mathbb N)}}(j)=A_{\operatorname{Add}_2}(j)=\infty$; a
topological nonextension theorem at alternating one-sided tails; maximal
online commitment deficit $B(n)=n$; maximal prefix erasure both on all ternary
inputs and on genuine addition inputs,
$\Delta(n)=\Delta_{\operatorname{Add}_2}(n)=n$; linear input-scale
reduced-residual growth $R_{\mathrm{in}}(n)\ge n$; the six-state output-delay quotient of the Berstel kernel,
including pairwise separation of the six reduced residuals; and the sharp
Berstel binary cleanup theorem
$\max_{u\in\{0,1\}^L,\ u=\operatorname{trim}_{\mathrm{MSD}}(u)}\tau(u)
=\lfloor L/2\rfloor$, together with the genuine-addition iteration-depth theorem
$\max_{w\in \operatorname{Add}_2^{\mathrm{MSD}}(n)}\tau(w)=\lceil n/2\rceil$.
The finite-word non-subsequentiality argument is now a corollary of the exact
prefix-erasure theorem.  A result-by-result table in the introduction makes
explicit which qualitative facts are classical and which quantitative
statements, domains, directions, and output requirements are proved here
(revised manuscript, Table~1, p.~3).  The current statement numbers and
theorem-by-theorem prior-art cross-references are collected in
Table~\ref{tab:response-comparison} of this response.  To the best of our
knowledge, the exact quantitative formulas in those tables are not stated in
the cited sources.

\item \textbf{Ten-state complexity of the Berstel adder.}
The manuscript no longer presents the existence of a ten-state Berstel adder
as a new contribution.  The abstract and introduction state that the
ten-state presentation is classical.  We also removed the earlier claim that
the induced unrestricted subsequential function has state complexity $10$.
The revised statement is: the ten raw complete-output residuals of the
displayed kernel are distinct, but after output-delay normalization the
induced subsequential function has the six-state quotient displayed in
Section 6, and the six reduced residuals are pairwise distinct.  The text now
cites Mohri's string-to-string minimization result (1994, revised Ref.~12)
and gives the direct normalized-residual lower bound before drawing the
minimality conclusion (Propositions~6.4--6.5, pp.~17--19).  The six-state
minimal quotient allows an initial output word; a formalism that forbids an
initial output uses one additional dummy initial state.

\item \textbf{Berstel adder contains the normalizer.}
A new remark in Section 5 explains the operational relationship: the Berstel
adder maps a ternary digitwise sum to a binary representation of the same
value, but the output need not be greedy; a further normalization step may be
needed to obtain the canonical Zeckendorf representative.  The new Section 7
then quantifies this operation: if $\tau(w)$ is the number of Berstel
iterations needed to reach the greedy representative, then on length-$n$
genuine addition inputs the worst case is exactly $\lceil n/2\rceil$.  The
upper bound is formulated through an explicit cascade-depth invariant
$D(u)$, measuring factors $(10)^r11$, and a direct argument from the Berstel
transition table for the first pass on the genuine addition language.  We also prove the
binary cleanup worst case $\lfloor L/2\rfloor$ and strengthen the lower bound
so that the extremal inputs are trimmed rather than leading-zero padded
(Theorems~7.5 and 7.9, pp.~23 and 26).

\item \textbf{Statement of the classical Berstel correctness result.}
The former heading ``Classical correctness'' has been replaced by ``Value
preservation of the Berstel recoding''.  The proof now explains that this
classical result is restated only to fix the transition, terminal-output, and
value conventions used in Sections~6--7; no novelty is claimed for the
correctness statement (Proposition~5.1, p.~15).  Its Fibonacci finite-word
source is now cited as Frougny (1999), Corollary 4 (p.~100; revised Ref.~8),
rather than Proposition 13 alone.

\item \textbf{Undefined terminology.}
The abstract and introduction now define or explain the terms that previously
appeared too early: LSDF/MSDF, the one-sided model, canonical normalizer,
coordinatewise unbounded anticipation, one-pass online delay, deterministic
subsequential transducer, and prefix-destruction index.  The abstract now
describes the latter as the largest number of initial normalized-output digits
that may be invalidated by extending an input.  Section 4 begins with
an explicit description of the LSDF one-sided model.  Definition~4.1 is our
adaptation of one-pass delay to the LSDF finite-support map; Frougny's 1999
paper uses a different MSDF online convention (Definition~4.1 and the
following remark, p.~8).  We also now define
$\mathcal A^{(\mathbb N)}$ for an alphabet with distinguished zero, identify
the all-zero sequence, and specify that LSDF trimming sends it to
$\varepsilon$.

\item \textbf{Notation.}
The notation $\oplus$ has been replaced by $\operatorname{Add}_{Z}$ for
canonical Zeckendorf addition.  The notation $\langle u\rangle$ has been
replaced by $\operatorname{enc}_{\mathrm{LSDF}}(u)$ to avoid conflict with
inner-product notation.  The alphabets in the online-delay definition are now
denoted by $\mathcal A$ and $\mathcal B$ rather than by
$\mathcal A_{\mathrm{in}}$ and $\mathcal A_{\mathrm{out}}$
(Definition~4.1, p.~8).

\item \textbf{Terminology such as overlay, stream, and low-to-high.}
The revised manuscript uses ``digitwise sum'', ``finite-support sequence'',
``LSDF'', and ``MSDF'' consistently.  The section title formerly using
``Overlay'' has been changed to ``Digitwise Sums and Normalization''
(Section~3, pp.~6--7).

\item \textbf{Value map and standard basis vector.}
The value map is now defined once on the full monoid $\mathcal D$, with the
restriction to admissible representatives stated explicitly.  The vector
$e_k$ is now defined as the finite-support sequence with digit $1$ at position
$k$ and digit $0$ elsewhere.  The canonical section is now typed explicitly
as $s:\mathcal D/\!\equiv\to\mathcal D$, with $q\circ s=\mathrm{id}$ and
$\operatorname{Fold}_\infty=s\circ q$ (Corollary~2.5, p.~5).

\item \textbf{Theorem statement and surrounding prose.}
The explanatory sentence following the addition theorem has been moved out of
the theorem statement.  Section 4 and its first subsection now start with
orientation sentences explaining the model and purpose of the section
(pp.~7--8).

\item \textbf{Iterates and local maps.}
The definition of $\Phi^t$ in the local lower bound now explicitly says that
it is the composition of $t$ copies of $\Phi$, and the radius is required to
satisfy $r\ge1$.  The conclusion is correspondingly restricted to a
time-homogeneous synchronous scheme obtained by iterating one fixed local rule
(Theorem~A.1, p.~27).

\item \textbf{Formal corrections in the quantitative invariants.}
The prefix-destruction indices now use
$n-|\operatorname{lcp}(\mathcal N(u),\mathcal N(v))|$, rather than attempting
to subtract a word from an integer.  We also replaced the output-truncated
residual count by the finite input-scale quantity
$R_{\mathrm{in}}(n)=\#\{\rho_u:|u|\le n\}$ and retained the explicit linear
lower bound $R_{\mathrm{in}}(n)\ge n$ (Theorems~4.6, 4.8, and 4.9,
pp.~11--13).

\item \textbf{Rewrite terminology and confluence.}
The proposition asserting preservation of value by the Fibonacci congruence
is now titled ``Value invariance'', rather than ``Completeness''.  The rewrite
system is terminating with a unique normal form, so the revised text now
states confluence as a corollary instead of disclaiming it
(Proposition~2.3, p.~5, and Corollary~3.2, p.~6).

\item \textbf{Berstel iteration results added in revision.}
In response to the observation that the Berstel adder contains a normalization
step but may need to be applied further, Section 7 now proves two exact bounds:
the binary cleanup depth is $\lfloor L/2\rfloor$, and the depth on length-$n$
genuine addition inputs is $\lceil n/2\rceil$.  The latter maximum is attained
by trimmed inputs and is therefore not a leading-zero padding artifact
(Theorems~7.5 and 7.9, pp.~23 and 26).
\end{enumerate}

\end{document}
